\documentclass{report}
\usepackage{amsfonts, amsmath}
\newcommand\R{{\mathbb R}}
\newcommand\Q{{\mathbb Q}}
\newcommand\N{{\mathbb N}}
\newcommand{\vf}{\varphi}
\pagestyle{empty}
\begin{document}
\hfuzz 12pt
\large
\centerline{\Huge \bf Analisi Matematica I}
\vskip.2cm
\centerline{\Huge Correzione Terzo Appello}

\vskip1cm
\begin{enumerate}
\item $1000000^{200}=(10^6)^{200}=10^{1200}$; $(1000!)^2\leq
(1000^{1000})^2 =10^{6000}$; $10!=3628800$, $2^4>10$ quindi
$2^{10!}\geq 10^{907200}$. D'altro canto $(1000!)^2\geq
(500^{500})^2\geq
(20^2)^{1000}\geq 10^{2000}$. Dunque
\[
1000000^{200}<(1000!)^2<2^{10!}.
\]
\item 
\[
\frac{q+1}{q-1}=1+\frac {2}{q-1}.
\]
Dunque, il numero di cui sopra puo essere intero se e solo se
$2(q-1)^{-1}\geq 1$ cio\`e $q\leq 3$. Da cui segue che gli unici
numeri permessi sono 2 e 3.
\item 
\[
\sum_{n=q+1}^\infty\frac 1{n!}\leq \frac 1{(q+1)!}\sum_{n=0}^\infty
\frac 1{(q+1)^n}\leq \frac 1{(q+1)!}\frac{1}{1-\frac 1{q+1}}
\leq \frac 1{q!q}.
\]
Veniamo alla seconda disuguaglianza.
\[
\frac pq-\sum_{n=0}^{q}\frac 1{n!}=\frac {p(q-1)!-\sum_{n=0}^{q}\frac
{q!}{n!}}{q!}
\]
dunque la sola possibilit\`a per la disuguaglianza di essere falsa \`e
se 
\begin{equation}\label{eq:abs}
p(q-1)!-\sum_{n=0}^{q}\frac
{q!}{n!}=0,
\end{equation}
dividendo il tutto per $(q-1)$ si ottiene
\[
p(q-2)!-\sum_{n=0}^{q-2}\frac{(q-2)!q}{n!}=\frac q{q-1}+\frac 1{q-1}=
\frac {q+1}{q-1}
\]
ma il menbro di sinistra \`e palesemente un numero intero mentre
quello di destra non lo \`e mai (vedi esercizio 2) dunque abbiamo una
contraddizione. Di consequenza \eqref{eq:abs} deve essere falsa,
cio\`e la disuguaglianza voluta deve essere vera. 

A questo punto si puo concludere: se $q>3$ da quanto visto segue che
\[
\left|e-\frac pq \right|\geq \left|\frac pq-\sum_{n=0}^q\frac
1{n!}\right|-\left| \sum_{n=q+1}^\infty\frac 1{n!}\right|\geq
\frac 1{q!}-\frac 1{q!q}>0
\]
Dunque $e$ non puo mai essere uguale ad un numero razionale a meno che
il denominatore non sia inferiore a $3$, ma 
\[
|e-\frac 83|=|e- (1+1+\frac 12+\frac 16)|\leq \frac 1{3!3}=\frac 1{18}
\]
quindi $e$ non puo essere ne $3$ ne $8/3$ ne altro numero razionale
con denominatore minore di $4$, e quindi non puo essere alcun numero
razione, cio\`e $e$ \`e un numero irrazionale. (Meditate gente, meditate)
\item Usando il criterio della radice si ha
\[
\lim_{n\to\infty}(1-1/n)^n=e^{-1}<1
\]
dunque la serie e' convergente.
Razionalizzando si ottiene
\[
\lim_{n\to\infty}\frac{n}{\sqrt{n^2+1}+n}=\frac 12.
\]
\item Il dominio \`e $x\in[-2,2]$ poiche per questi valori la
quantita' sotto ambedue le radici \`e positivo. Poi il grafico \`e
simmetrico attorno ad $0$. In pi\`u in zero la funzione non \`e
differenziabile. Per vederlo si espanda la prima radice usando Taylor
al secondo ordine. Si ottiene
\[
\sqrt{2-(2+(x/2)^2+{\mathcal O}(x^4))}=\frac{|x|}2+{\mathcal
O}(|x|^3).
\] 
Inoltre \`e facile vedere che per ogni $x> 0$ la derivata \`e
positiva. Dunque la funzione ha un minimo in zero, ma in zero ha uno
spigolo. 
\item Prima di tutto, \`e facile vedere dal grafico dell'esercizio
precedente che $f(x)\leq\sqrt2 x$, per $x>0$. Da questo segue
\[
\ell_{n}=f(\ell_{n-1})\leq(\sqrt 2)^{-1}\ell_{n-1}\leq (\sqrt 2)^{-n}\ell_0=
2(\sqrt 2)^{-n} .
\]
Dunque $\lim_{n\to\infty}\ell_n=0$ (e questo, come ho detto durante lo
scritto, assicurava gia almeno 4 punti).
Tuttavia per rispondere alla domanda \`e necessaria una stima pi\`u
precisa. Raffiniamo la stima gia visto nell'esercizio precedente: per
$x>0$
\[
f(x)= \frac x2 +\frac {x^3}{32}+{\mathcal O}(x^5)
\]
dove abbiamo usato Taylor al terzo ordine per sviluppare la radice.
A questo punto, se $x$ \`e sufficientemente piccolo, si ha
\[
f(x)\leq \frac x2(1+\frac {x^2}8)\leq \frac x2 e^{\frac{x^2}8}.
\]
Ora, poiche $\ell_n$ tende a zero, eriste un $n_0$ tale che per tutti
gli $n>n_0$ $\ell_n$ \`e abbastanza piccolo nel senso di cui sopra.

Usiamo ora questa disuguaglianza per iterare, come abbiamo fatto sopra,
\[
\begin{split}
\ell_n=&f(\ell_{n-1})\leq \ell_{n-k}2^{-k}e^{\frac
18\sum_{j=1}^k\ell_{n-j}^2}\leq \ell_{n_0}2^{-n+n_0}e^{\frac
14\sum_{j=1}^{n-n_0}\ell_{n-j}^2}\\
&\leq \ell_{n_0}2^{-n+n_0}e^{\frac
18\sum_{j=1}^{\infty}\ell_{j}^2}.
\end{split}
\]
Dunque, usando la stima gia ottenuta 
\[
\ell_n\leq \ell_{n_0}2^{-n+n_0}e^{\frac
12\sum_{j=1}^{\infty}2^{-j}}\leq \ell_{n_0}2^{-n+n_0}e.
\]
Finalmente,
\[
2^n\ell_n\leq 2^{n_0}\ell_{n_0}e
\]
che non dipende da $n$ e dunque mostra che la succesione $a_n$ \`e
limitata.

Per i {\bf curiosi}, il {\sl segreto} di questo esercizio \`e il seguente.
Un piccolo calcolo trigonomerico (o un semplice disegno) mostra che,
per ogni$\vf\in[0,\pi/2]$,
\[
f(2\sin\vf)=2\sin\frac\vf 2.
\]
Allora se si pone $\vf_0=\pi$ e $\vf_{n+1}=\frac{\vf_n}2$ si ha
$\ell_n=2\sin\vf_{n+1}$. Il semplice disegno a cui facevo riferimento
sopra consiste nel realizzare che $\ell_n$ \`e il lato del poligono
regolare con $2^{n+1}$ lati. Dunque $a_n$ non \`e altro che il
semiperimentro di tale poligono. Poich\`e il perimetro dei poligoni
inscritti dovrebbe convergere alla lunghezza della circoferenza
(Archimede docet) si deve avere che
\[
\lim_{n\to\infty}a_n=\pi,
\]
dunque, non solo la succesione \`e limitata ma converge pure e a
nientedimeno che al nostro
amico $\pi$.
\end{enumerate}
\end{document}
